\subsection{Genuine coefficient extension and domination of line components}
\label{subsec:scalar-extended-line-components}

\begin{theorem}[Whole tensor comparison for polynomial principal opens]
\label{algebra:principal-open-scalar-tensor}
\lean{EnvelopeProperty.AffineLine.principalOpenTensorEquiv,
EnvelopeProperty.AffineLine.principalOpenTensorEquiv_tmul}
\leanok
For arbitrary commutative rings $k,F$ with a genuine coefficient map and
any polynomial $g\in k[X_i]$, the actual ring isomorphism
\[
 F\otimes_k k[X_i,1/g]\simeq
 F[X_i,1/g_F]
\]
is an $F$-algebra isomorphism, where $g_F$ is the entire coefficient-mapped
polynomial. It preserves every original tensor $s\otimes q$ as the
localization of $s q_F$. No nonzero or injectivity hypothesis is needed.
\end{theorem}
\begin{proof}
Actual tensor localization first identifies the original tensor ring with
the localization of $F\otimes_k k[X_i]$ at $1\otimes g$. The genuine
polynomial tensor isomorphism maps every $s\otimes q$ to $s q_F$ and
maps the entire denominator $1\otimes g$ to $g_F$. It therefore maps
the actual power submonoid onto that of $g_F$. The true localization
isomorphism induced by this full polynomial isomorphism gives the
displayed entire ring isomorphism. Its genuine tensor and localization
evaluation formulas prove the assertion on all original polynomials
and scalars, including zero-localization cases.
\end{proof}

\begin{theorem}[Actual scalar-extended ambient map retains finiteness and injectivity]
\label{algebra:principal-open-scalar-map}
\lean{EnvelopeProperty.AffineLine.principalOpenScalarMap,
EnvelopeProperty.AffineLine.principalOpenScalarMap_tensor,
EnvelopeProperty.AffineLine.principalOpenScalarMap_algebraMap,
EnvelopeProperty.AffineLine.principalOpenScalarMap_finite,
EnvelopeProperty.AffineLine.principalOpenScalarMap_injective}
\leanok
\uses{algebra:principal-open-scalar-tensor,spectrum:scalar-extension-algebra-cancellation}
Every genuine $k$-algebra homomorphism
$\varphi:k[X_i,1/g]\to B$ gives an actual $F$-algebra homomorphism
\[
 \varphi_F:F[X_i,1/g_F]\longrightarrow F\otimes_k B.
\]
Its full diagram agrees with $1\otimes\varphi$ on every entire tensor
and every original polynomial. A finite original map gives a finite
new map. If $F$ is a genuinely flat $k$-module and the original map is
injective, then the new whole map is injective. This map construction
and its finiteness statement use arbitrary commutative rings.
\end{theorem}
\begin{proof}
Compose the inverse of the preceding actual localized tensor isomorphism
with the genuine tensor algebra map $1\otimes\varphi$. Its inverse law
gives the whole diagram on every tensor, not only generators. The proved
entire polynomial tensor formula identifies the original polynomial
inclusion with its exact coefficient-mapped inclusion in this diagram.

For finiteness, equip $B$ with the actual action induced by $\varphi$.
The original coefficient commutation gives its real scalar tower.
Actual scalar tensor cancellation identifies $F\otimes_k B$ with the
base change of this original finite module to
$F\otimes_k k[X_i,1/g]$, retaining the full base-ring action. Genuine
module finiteness after this actual cancellation, followed by the
localization isomorphism, yields finiteness of the new entire ring map.
For injectivity, true coefficient flatness preserves the injectivity
of $\varphi$ as an actual linear map under left tensoring. Composing
with the preceding inverse ring isomorphism preserves injectivity.
No free or finite-flat hypothesis on the original algebra is used.
\end{proof}

\begin{theorem}[Actual scalar-extended minimal components dominate the line]
\label{algebra:scalar-line-component-domination}
\lean{EnvelopeProperty.AffineLine.scalarComponentLineHom_injective,
EnvelopeProperty.AffineLine.scalarComponent_quotient_ringKrullDim_eq_one,
EnvelopeProperty.AffineLine.scalarComponentLineHom_comap_surjective}
\leanok
\uses{algebra:geometric-domain-base-change,algebra:principal-open-scalar-map,
algebra:line-component-domination}
Let $k$ be algebraically closed, $F/k$ any genuine field extension,
and $B$ an original domain $k$-algebra. Let the actual ambient map
$k[X_i,1/g]\to B$ be finite and injective, with $g\ne0$ and finitely
many coordinates. Choose a real line over $F$ with a nonzero parameter
coordinate and $g_F(a+bT)\ne0$. For every actual minimal prime $P$ over
the original transverse-equation image under the constructed
$\varphi_F$, the genuine whole component-line map is injective,
\[
 \dim((F\otimes_k B)/P)=1,
\]
and its prime-spectrum map onto the original open line is surjective.
Its finiteness and full scheme square are the preceding actual component
producers. Domain of $F\otimes_k B$ is a conclusion of a real algebraic
producer, not an input certificate in these theorems.
\end{theorem}
\begin{proof}
Genuine geometric domain preservation over the algebraically closed
original field proves that the entire tensor ring $F\otimes_k B$ is a
domain. The actual field coefficient map is injective, so its entire
polynomial map preserves $g\ne0$. A field module is flat over $k$;
the preceding actual map producers therefore make $\varphi_F$ finite
and injective. They also preserve all original coefficient and tensor
elements.

Apply the proved real component-line construction to this actual new
map and the supplied original minimal prime of its transverse-equation
image. Its genuine height and dimension arguments derive whole line-map
injectivity and quotient dimension one. Actual integral lying-over
gives surjectivity on the entire prime spectrum. The ambient tensor-domain
and finite-map hypotheses used there have now been discharged by genuine
coefficient-extension producers, rather than encoded in a component or
geometric-integrality record. No original normality, completeness,
separability, infinitude or characteristic hypothesis is added.
\end{proof}

This gives the scalar-extended ambient-domain and line-component domination
input of FHX Proposition~5.22 from a real original finite injective affine
map. Existence of the paper's complete algebraically closed valued field
and rational lift, selection of a line avoiding the bad open complement,
finite normalization and its smooth or flat structure, and the curve
potential argument remain separate genuine constructions.
